%==============================================================================
% Externalities and Market Inefficiency - Endogenous Growth
% Arrow Learning-by-Doing Framework
%==============================================================================

\subsection*{Arrow Learning-by-Doing Framework}

\medskip
\noindent\textbf{Individual Firm Production:}

\begin{align*}
Y_j = \bar{A} K_j^\alpha L_j^{1-\alpha} \tag{Firm Production}
\end{align*}

\medskip
\noindent\textbf{External Effect (Learning-by-Doing):}

\begin{align*}
\bar{A} = A_0 \left(\int K_i \, di\right)^\eta = A_0 K^\eta, \qquad \eta > 0
\tag{Externality}
\end{align*}

\medskip
\noindent\textbf{Aggregate Production:}

\begin{align*}
Y = A_0 K^{\alpha+\eta} L^{1-\alpha} \equiv A K^{\alpha+\eta} \tag{Aggregate}
\end{align*}

\medskip
\begin{tcolorbox}[colback=blue!5,colframe=blue3,title=Key Insight]
Firms treat $\bar{A}$ as given, but each firm's capital contributes to aggregate knowledge, raising $\bar{A}$ for all.
\end{tcolorbox}

\clearpage

%==============================================================================
\subsection*{The Knife-Edge Case (Solow-AK)}

\begin{derivationbox}
Set $\alpha + \eta = 1$ (externality exactly offsets diminishing returns):

\begin{align*}
Y &= A_0 K^{\alpha+\eta} L^{1-\alpha} = A_0 K L^{1-\alpha} \\
  &= A K \quad \text{where } A = A_0 L^{1-\alpha}
\end{align*}

\medskip
Output is linear in capital. This gives a \textbf{constant marginal product of capital}.
\end{derivationbox}

\begin{resultbox}
\begin{align*}
\boxed{Y = AK} \tag{AK Production}
\end{align*}
\end{resultbox}

\medskip
\noindent Under AK production:
\[
\text{MPK} = \frac{\partial Y}{\partial K} = A = \text{constant}
\]

\clearpage

%==============================================================================
\subsection*{Private vs. Social Marginal Product}

\begin{derivationbox}
\textbf{Private MPK} (firm takes $A$ as given):
\[
r + \delta = \frac{\partial Y_j}{\partial K_j} = \alpha \bar{A} K_j^{\alpha-1} L_j^{1-\alpha} = \alpha A
\tag{Private MPK}
\]

\medskip
\textbf{Social MPK} (accounts for externality):
\[
r + \delta = \frac{\partial Y}{\partial K} = (\alpha + \eta) A
\tag{Social MPK}
\]
\end{derivationbox}

\begin{resultbox}
\[
\boxed{r + \delta = \underbrace{\alpha A}_{\text{Private MPK}} < \underbrace{(\alpha + \eta)A}_{\text{Social MPK}}}
\]
\end{resultbox}

\medskip
\begin{keyformulabox}
Since $\eta > 0$: \textbf{Social MPK > Private MPK}
\end{keyformulabox}

\clearpage

%==============================================================================
\subsection*{Growth Comparison: Private vs. Planner}

\medskip
\noindent With CRRA utility $u(c) = \frac{c^{1-\sigma} - 1}{1-\sigma}$ and AK technology:

\medskip
\noindent\textbf{Private Economy Growth Rate:}
\begin{align*}
\gamma_c = \frac{\alpha A - \delta - \rho}{\sigma}
\tag{Private Growth}
\end{align*}

\medskip
\noindent\textbf{Social Planner Growth Rate:}
\begin{align*}
\gamma_c^{SP} = \frac{(\alpha + \eta)A - \delta - \rho}{\sigma}
\tag{Planner Growth}
\end{align*}

\medskip
\begin{resultbox}
\[
\boxed{\gamma_c^{SP} = \gamma_c + \frac{\eta A}{\sigma} > \gamma_c}
\]
\end{resultbox}

\medskip
\begin{tcolorbox}[colback=orange!10,colframe=orange3,title=Result]
The planner grows faster because it internalizes the positive externality from learning-by-doing.
\end{tcolorbox}

\clearpage

%==============================================================================
\subsection*{Inefficiency Result and Policy}

\begin{derivationbox}
The market economy is inefficient because:

\begin{enumerate}
    \item \textbf{Positive externality}: Each firm's investment increases $A$ for all firms
    \item \textbf{Under-accumulation}: Private agents don't internalize the learning-by-doing effect
    \item \textbf{Result}: Market saves less than the socially optimal level
\end{enumerate}
\end{derivationbox}

\begin{resultbox}
\begin{center}
\textbf{Social MPK > Private MPK} $\Rightarrow$ \textbf{Market under-accumulates capital}
\end{center}
\end{resultbox}

\medskip
\subsection*{Policy Implication}

\begin{center}
\begin{tabular}{ll}
\toprule
\textbf{Problem} & \textbf{Solution} \\
\midrule
Positive externality & Capital accumulation subsidy \\
from learning-by-doing & to align private and social returns \\
\bottomrule
\end{tabular}
\end{center}

\medskip
The optimal subsidy raises the private return to match the social return, inducing higher savings and growth.

\clearpage

%==============================================================================
\subsection*{Summary Table: Private vs. Social Planner}

\begin{center}
\begin{tabular}{lccc}
\toprule
 & Private Economy & Social Planner & Difference \\
\midrule
Production Function & $Y = A_0 K^\alpha L^{1-\alpha}$ & $Y = A_0 K^{\alpha+\eta} L^{1-\alpha}$ & $\eta$ term \\
Effective Returns & $\alpha A$ & $(\alpha + \eta)A$ & $+\eta A$ \\
Growth Rate & $\frac{\alpha A - \delta - \rho}{\sigma}$ & $\frac{(\alpha+\eta)A - \delta - \rho}{\sigma}$ & $+\frac{\eta A}{\sigma}$ \\
Capital Level & Lower & Higher & Subsidy needed \\
\bottomrule
\end{tabular}
\end{center}

\medskip
\begin{tcolorbox}[colback=green!10,colframe=green3,title=Key Takeaway]
The positive learning-by-doing externality ($\eta > 0$) creates a wedge between private and social returns, leading to under-accumulation of capital. A capital subsidy can restore efficiency.
\end{tcolorbox}
